\section{Action 1.5 : Rattachement des Projets à la Solution (dotnet sln add)}

Après avoir initialisé indépendamment les trois briques du système (\texttt{Shared}, \texttt{Back} et \texttt{Front}), l'étape d'ingénierie consiste à les fédérer officiellement sous l'autorité de la solution racine \textbf{\texttt{FullStackCSharp.sln}}.

Cette opération s'effectue via la commande en ligne de commande :
\begin{center}
\texttt{dotnet sln FullStackCSharp.sln add FullStackCSharp.Shared FullStackCSharp.Back FullStackCSharp.Front}
\end{center}

Loin d'être une simple commodité visuelle pour les environnements de développement, ce rattachement remplit trois fonctions capitales dans le cycle de vie logiciel :
\begin{itemize}
    \item \textbf{Indexation officielle des chemins relatifs} : La solution consigne le chemin de chaque fichier \texttt{.csproj} par rapport à la racine, permettant à tout développeur ou serveur d'intégration continue de cloner le dépôt et de compiler l'ensemble sans reconfiguration.
    \item \textbf{Attribution des identifiants uniques (GUIDs)} : MSBuild assigne à chaque projet un identifiant machine immuable garantissant l'intégrité du graphe de dépendances.
    \item \textbf{Unification de la matrice de build} : La solution projette les configurations globales (\texttt{Debug} et \texttt{Release}) ainsi que les architectures cibles sur chaque projet membre.
\end{itemize}

\vspace{0.3cm}

\begin{schemabox}{Topologie de Coordination par la Solution Fédératrice}
\begin{center}
\begin{tikzpicture}[
    boxslnhub/.style={
        rectangle, 
        draw=accent!90!black, 
        fill=accent!7!white, 
        rounded corners=2.5mm, 
        minimum width=15.4cm, 
        minimum height=1.2cm, 
        inner ysep=2.5mm,
        align=center, 
        line width=1.1pt
    },
    boxproj/.style={
        rectangle, 
        draw=#1!80!black, 
        fill=#1!6!white, 
        rounded corners=2mm, 
        minimum width=4.7cm, 
        minimum height=1.9cm, 
        inner ysep=3mm,
        align=center, 
        line width=0.85pt
    },
    boxmatrix/.style={
        rectangle, 
        draw=primary, 
        fill=primary!95!black, 
        rounded corners=2.5mm, 
        minimum width=15.4cm, 
        minimum height=1.2cm, 
        inner ysep=2.5mm,
        align=center, 
        line width=1pt
    },
    brancharrow/.style={
        -{Stealth[length=2.8mm]}, 
        thick, 
        color=accent!90!black
    },
    matrixarrow/.style={
        -{Stealth[length=2.8mm]}, 
        thick, 
        color=primary!80!white
    }
]

% Top: Solution Root Hub (y = 3.6)
\node[boxslnhub] (sln) at (0, 3.6) {
    \textbf{\normalsize Orchestrateur Racine : \texttt{FullStackCSharp.sln}}\\[2pt]
    \footnotesize Indexation textuelle des projets $\bullet$ Fédérateur de compilation $\bullet$ Point d'entrée outillage
};

% Middle Row: 3 Federated Projects (y = 1.0)
\node[boxproj=cyanproj] (shared) at (-5.2, 1.0) {
    \textbf{\normalsize FullStackCSharp.Shared}\\[2pt]
    {\color{cyanproj}\rule{3.8cm}{0.4pt}}\\[2pt]
    \scriptsize Type : \texttt{classlib} (.NET 10 LTS)\\[1pt]
    \tiny GUID : \texttt{\{3F385EC1-20BC-4DD2...\}}\\
    \tiny Rôle : Contrats \& DTOs
};

\node[boxproj=secondary] (back) at (0, 1.0) {
    \textbf{\normalsize FullStackCSharp.Back}\\[2pt]
    {\color{secondary}\rule{3.8cm}{0.4pt}}\\[2pt]
    \scriptsize Type : \texttt{webapi} (.NET 10 LTS)\\[1pt]
    \tiny GUID : \texttt{\{326CA162-2E67-430D...\}}\\
    \tiny Rôle : Kestrel \& SQLite
};

\node[boxproj=indigo] (front) at (5.2, 1.0) {
    \textbf{\normalsize FullStackCSharp.Front}\\[2pt]
    {\color{indigo}\rule{3.8cm}{0.4pt}}\\[2pt]
    \scriptsize Type : \texttt{maui} (.NET 10 LTS)\\[1pt]
    \tiny GUID : \texttt{\{D25ACDE3-554B-495E...\}}\\
    \tiny Rôle : UI Cross-Plateforme
};

% Branching arrows from SLN to projects with clean labels
\draw[brancharrow] (-5.2, 3.0) -- (shared.north)
    node[midway, left=4pt, font=\tiny\bfseries\color{accent!90!black}] {add project};

\draw[brancharrow] (0, 3.0) -- (back.north)
    node[midway, right=4pt, font=\tiny\bfseries\color{accent!90!black}] {add project};

\draw[brancharrow] (5.2, 3.0) -- (front.north)
    node[midway, right=4pt, font=\tiny\bfseries\color{accent!90!black}] {add project};

% Bottom: Matrix of configurations (y = -1.5)
\node[boxmatrix] (matrix) at (0, -1.5) {
    \color{white}\textbf{\normalsize Matrice Globale de Compilation MSBuild (Debug | Release)}\\[2pt]
    \footnotesize\color{emerald!30!white}\texttt{GlobalSection(SolutionConfigurationPlatforms)} $\bullet$ \color{cyanproj!30!white}\texttt{Any CPU / x64} $\bullet$ \color{accent!30!white}Synchronisation unifiée
};

% Connecting arrows from projects to matrix
\draw[matrixarrow] (shared.south) -- (-5.2, -0.9);
\draw[matrixarrow] (back.south) -- (0, -0.9);
\draw[matrixarrow] (front.south) -- (5.2, -0.9);

\end{tikzpicture}
\end{center}
\end{schemabox}

\vspace{0.2cm}

\begin{bestpracticebox}{L'Organisation en Dossiers Virtuels (Solution Folders)}
Dans les solutions d'entreprise comportant des dizaines de projets (APIs, microservices, tests unitaires, benchmarks), il est recommandé d'organiser les projets dans des \textbf{Solution Folders} virtuels (ex : \texttt{src/}, \texttt{tests/}, \texttt{docs/}).
Ces dossiers logiques sont créés directement via la CLI (\texttt{dotnet sln add ... --solution-folder src}) et structurent l'arborescence dans l'éditeur sans modifier l'emplacement physique des répertoires sur le disque.
\end{bestpracticebox}

\newpage

\subsection*{Anatomie d'une Entrée Projet dans le Fichier .sln}

Lorsque la commande \texttt{dotnet sln add} est exécutée, elle modifie le fichier texte de la solution en inscrivant deux blocs de données essentiels :

\begin{description}
    \item[\color{secondary}1. Déclaration du Projet (\texttt{Project})] : Définit le type de projet via son GUID de langage (ex : \texttt{\{FAE04EC0-...\}} pour le C\#), le nom logique du composant, son chemin relatif sur le disque, et lui attribue son propre GUID de solution.
    \item[\color{cyanproj!90!black}2. Mappage des Configurations (\texttt{ProjectConfigurationPlatforms})] : Associe chaque projet aux configurations de la solution (\texttt{Debug|Any CPU}, \texttt{Release|Any CPU}). Cela indique à MSBuild s'il doit inclure ou ignorer le projet lors de la compilation de la solution dans un mode particulier.
\end{description}

\vspace{0.3cm}

\subsection*{Questions Fréquentes du Débutant}

\faqitem{Que fait concrètement la commande \texttt{dotnet sln add} dans le fichier \texttt{.sln} ?}{
    Elle n'altère absolument pas le code de vos projets (\texttt{.csproj} ou fichiers \texttt{.cs}). Elle ajoute simplement quelques lignes de texte descriptif dans le fichier \texttt{FullStackCSharp.sln} pour enregistrer le chemin relatif du projet et déclarer qu'il participe désormais aux compilations globales déclenchées depuis la racine.
}

\faqitem{Si un projet n'est pas rattaché au fichier \texttt{.sln}, peut-on quand même le compiler ?}{
    \textbf{Oui, absolument.} Chaque projet \texttt{.csproj} est une entité de compilation autonome. Vous pouvez vous placer dans le sous-dossier \texttt{FullStackCSharp.Back} et exécuter \texttt{dotnet build} ou \texttt{dotnet run} avec succès. Le rattachement au \texttt{.sln} devient indispensable lorsque vous souhaitez compiler ou tester l'intégralité de la solution en une seule commande depuis la racine.
}

\faqitem{Est-ce que rattacher les projets à la solution crée les dépendances (Back $\rightarrow$ Shared, Front $\rightarrow$ Shared) ?}{
    \textbf{Non, absolument pas ! C'est une confusion fréquente chez les développeurs débutants.}
    La commande \texttt{dotnet sln add} regroupe les projets au sein d'un même espace de travail pour que MSBuild puisse les compiler ensemble en une seule commande. Cependant, elle \textbf{n'ajoute aucune référence de code ni liaison binaire} entre les projets.
    Pour que le projet \texttt{Back} ou \texttt{Front} puisse consommer les DTOs de \texttt{Shared}, il faut explicitement ajouter une balise \texttt{<ProjectReference>} dans les fichiers \texttt{.csproj} respectifs (via \texttt{dotnet add reference}). Ce câblage précis fait l'objet de l'\textbf{Action 1.7}.
}

\faqitem{Pourquoi le fichier \texttt{.sln} utilise-t-il des GUIDs hexadécimaux pour identifier les projets ?}{
    Les GUIDs (\textit{Globally Unique Identifiers}) garantissent que chaque projet possède une identité mathématiquement unique et immuable dans tout le système, même si deux projets portent le même nom dans des répertoires différents ou si un dossier est renommé sur le disque. MSBuild s'appuie sur ces clés pour relier sans ambiguïté les configurations de compilation.
}

